\documentclass[11pt,a4paper]{article}

\usepackage[brazil]{babel}
\usepackage[utf8]{inputenc}
\usepackage[T1]{fontenc}
\usepackage[left=1.5cm,right=1.5cm,top=1.5cm,bottom=1.5cm]{geometry}
\usepackage{booktabs}
\usepackage{multirow}
\usepackage{longtable}
\usepackage{array}
\usepackage{pdflscape}
\usepackage{tabularx}
\usepackage[hidelinks]{hyperref}
\usepackage{caption}

\renewcommand{\arraystretch}{1.2}
\setlength{\LTpre}{0pt}
\setlength{\LTpost}{0pt}

\begin{document}

\section*{Modelo de Tabela Complexa}

Este arquivo demonstra o uso de células mescladas, colunas agrupadas, alinhamento refinado e estruturação de tabelas mais densas, úteis em relatórios técnicos e documentos editoriais.

\begin{landscape}
\small
\begin{longtable}{p{3.0cm} p{3.8cm} p{3.4cm} p{4.6cm}}
\caption{Quadro de organização editorial e técnica} \label{tab:quadro}\\
\toprule
\textbf{Categoria} & \textbf{Subprocesso} & \textbf{Critério} & \textbf{Observações} \\
\midrule
\endfirsthead

\toprule
\textbf{Categoria} & \textbf{Subprocesso} & \textbf{Critério} & \textbf{Observações} \\
\midrule
\endhead

\multirow{4}{*}{Pré-publicação}
& Submissão & Metadados completos & Título, autor, afiliação e resumo devidamente preenchidos. \\
& Triagem inicial & Conformidade formal & Verificação de norma, formatação e documentação obrigatória. \\
& Revisão textual & Clareza e padronização & Correção de estilo, ortografia e consistência visual. \\
& Controle de versão & Histórico de alterações & Registro de mudanças e organização dos arquivos. \\

\midrule
\multirow{4}{*}{Publicação}
& Copyediting & Ajustes de linguagem & Adequação da redação ao padrão editorial da revista. \\
& Diagramação & Estrutura visual & Organização de títulos, subtítulos, tabelas e figuras. \\
& Prova final & Conferência técnica & Última revisão antes da disponibilização ao público. \\
& Indexação & Publicação final & Inserção de dados, DOI e metadados do artigo. \\

\midrule
\multirow{3}{*}{Pós-publicação}
& Divulgação & Comunicação institucional & Publicação em canais oficiais e redes sociais. \\
& Monitoramento & Acompanhamento & Verificação de acesso, leitura e estabilidade dos arquivos. \\
& Arquivamento & Preservação documental & Organização da versão final e backup do material. \\

\bottomrule
\end{longtable}
\end{landscape}

\section*{Exemplo adicional com células mescladas}

\begin{table}[h]
\centering
\caption{Resumo de tarefas por etapa}
\begin{tabular}{|p{3.2cm}|p{3.8cm}|p{3.0cm}|}
\hline
\textbf{Etapa} & \textbf{Responsável} & \textbf{Entrega esperada} \\
\hline
\multirow{2}{*}{Revisão}
& Editor & Ajustes gerais do texto \\
\cline{2-3}
& Assistente editorial & Conferência de dados e estilo \\
\hline
\multirow{2}{*}{Produção}
& Diagramador & Organização visual do conteúdo \\
\cline{2-3}
& Autor & Validação final do material \\
\hline
\end{tabular}
\end{table}

\end{document}